\section{The distribution problem and its preliminary tools}\label{sec:distributions}

A \emph{unit} is an ordered pair of raw vertices with no hole between its endpoints.
The endpoints of a unit may be dependent.
When two units are sampled independently, their four endpoints need not be independent within either unit.
The event of interest is that all four cross pairs have holes, as in \Cref{fig:four-holes}.

The raw reference law \(\rawlaw\) is the uniform law of the exact frames in
\Cref{eq:raw-frame-law}.  We will replace its conditional channel law by a
thinned law \(\pi\), while preserving its primal marginal exactly.
The reference image and peeling arguments below use the absolute joint cap
\begin{equation}\label{eq:raw-law-caps}
 \sigma\le 2^{DN}\rawlaw^2.
\end{equation}
This condition alone imposes no constant bound on either endpoint marginal.
The stronger marginal hypotheses required for the distribution tests will
always be stated separately, relative to \(\pi\).

\subsection{Uniform frame laws}

We record the linear algebra behind all reference image laws.
For a subspace of a nominal coefficient space, its \emph{image} is its image under the relevant frame map.
Nominal independence is independence before applying these maps.

\begin{lemma}\label{lem:paired-frame-orbit}
Let \(A,A':U\to V\) and \(B,B':W\to V^*\) be injective maps between finite-dimensional spaces over \(\F\).
If
\[
 B\trans A=(B')\trans A',
\]
there is \(L\in\GL(V)\) with \(LA=A'\) and \(L^{-\mathsf T}B=B'\).
Thus based individually injective plus and minus frames with a prescribed Gram matrix form one orbit under the simultaneous primal and dual action.
\end{lemma}
\begin{proof}
Write \(\Phi=B\trans:V\to W^*\) and \(\Phi'=(B')\trans:V\to W^*\).
Both are surjective.
The based map \(f:A(U)\to A'(U)\) given by \(f(Au)=A'u\) satisfies \(\Phi'f=\Phi\).
It therefore maps \(A(U)\cap\ker\Phi\) isomorphically onto \(A'(U)\cap\ker\Phi'\).
Extend this restriction to an isomorphism \(L_0:\ker\Phi\to\ker\Phi'\).

Choose a complement \(U_0\) of \(A(U)\cap\ker\Phi\) in \(A(U)\).
Its images under \(\Phi\) are independent.
Extend a basis of \(\Phi(U_0)\) to a basis of \(W^*\), and choose lifts of the added basis vectors under \(\Phi\) and \(\Phi'\), respectively.
For the basis vectors in \(\Phi(U_0)\), choose the lifts in \(U_0\) and their images under \(f\).
Define \(L\) to equal \(L_0\) on \(\ker\Phi\) and to carry each chosen lift to its primed lift.
These are direct-sum decompositions of \(V\), so \(L\) is an isomorphism.
It extends \(f\) and satisfies \(\Phi'L=\Phi\), which is equivalent to the two required frame identities.
\end{proof}

\begin{lemma}\label{lem:frame-images}
The following statements hold for all sufficiently large \(n\).
\begin{enumerate}[label=\textup{(\alph*)}]
\item
At independently sampled raw vertices and components, images of prescribed primal coefficient tuples that are independent separately at each vertex and sign are uniform subject to their individual injectivities and their prescribed same-vertex Gram matrices.
For tuples using channels as well, the same assertion holds conditionally on the full self-channel Gram matrices \(X_e\trans Y_e\); the relevant full Gram matrix is
\(\begin{psmallmatrix}E&0\\0&X_e\trans Y_e\end{psmallmatrix}\).

\item
Let a fixed tuple of nominal directions use arbitrary combinations of the two endpoints' primal and channel coefficients.
Compute its rank separately in each component and sign in the \emph{direct sum over both endpoints}, and let the total be \(t\).
For every preassigned \(\epsilon>0\), choosing \(M_0\) sufficiently large gives, under \(\rawlaw^2\),
\begin{equation}\label{eq:reference-image-cap}
 \PP\{\text{all those images equal prescribed vectors}\}
 \le 2^{-(1-\epsilon)tN}\qquad(t\ge1).
\end{equation}
This is valid for all possible ranks \(t\), including \(t=O(n)\), not only for bounded tuples.

\item
The plus frame \([P,X]\) is uniform among injective maps.
Conditional on it, the minus columns are independent uniforms on the affine spaces prescribed in \Cref{eq:raw-frame-law}, followed only by the conditioning that \([Q,Y]\) is injective.
The failure probability of this last condition, before conditioning, is at most \(2^{2(d+h)-N}\) per component.

Given \(P,Q\), the channels \(X,Y\) are independent uniform-column draws on \(\ker Q\trans\) and \(\ker P\trans\), respectively, conditioned on \([P,X]\) and \([Q,Y]\) being injective.
The two rank-failure probabilities are exponentially small in \(N\).
Each channel frame alone is uniform among injective maps \(\F^h\to\F^N\).
For any fixed \(q\) independent vectors in its paired ambient space, the probability that transpose evaluation on their span is not injective is at most
\begin{equation}\label{eq:channel-transpose-failure}
 2^{q-h+1}.
\end{equation}
The unconditional joint channel law is bounded above by a constant, depending on \(g,h\) but not on \(n\), times the completely independent uniform-column law.

\item
For \(p\) independent uniform plus vectors and \(q\) independent uniform minus vectors in paired copies of \(\F^N\), and a prescribed \(p\)-by-\(q\) Gram matrix \(G\), the probability of that Gram matrix together with individual injectivity of both frames is at most \(2^{-pq}\) and at least
\begin{equation}\label{eq:bounded-gram-normalization}
 2^{-pq}(1-2^{p-N})(1-2^{p+q-N}).
\end{equation}
In particular, any consistent Gram specification on a bounded number of vectors, together with injectivity on each side, has probability bounded below by a positive constant for large \(N\).
\end{enumerate}
\end{lemma}
\begin{proof}
The raw law is invariant under \(L\) on each \(V_e^+\) and \(L^{-\mathsf T}\) on \(V_e^-\), independently at vertices and components.
For primal tuples, their restricted Gram is fixed by \(E\).
By \Cref{lem:paired-frame-orbit}, all based realizations with that Gram and individual injectivity lie in one orbit.
Projection of an invariant uniform raw law is uniform on this orbit: an element carrying one partial realization to another bijects their full-frame completions.
There is at least one completion because these tuples are restrictions of a raw frame.
After fixing \(X_e\trans Y_e\), the same argument applies to arbitrary full-frame coefficient tuples.
Independence over raw draws and components proves (a).

For (c), the action is transitive on injective plus frames, so their marginal is uniform.
Given a plus frame, each column of \(Q\) has specified evaluations against \(P,X\), while each column of \(Y\) has zero evaluations against \(P\).
There is no other constraint until full minus injectivity is imposed.
The remaining allowed set is therefore the product of those affine column spaces restricted to full rank.

Each column space contains the common translation space
\(\operatorname{ann}(\im[P,X])\), of dimension \(N-d-h\).
For a nonzero coefficient vector \(c\in\F^{d+h}\), the corresponding linear combination of minus columns is consequently uniform on an affine space containing that translation space.
Its chance of being zero is at most \(2^{-(N-d-h)}\).
Union over the fewer than \(2^{d+h}\) nonzero coefficient vectors gives the asserted bound \(2^{2(d+h)-N}\).

Alternatively, condition first on \(P,Q\).
The allowed \(X\)-columns lie in \(\ker Q\trans\), the allowed \(Y\)-columns in \(\ker P\trans\), and their two full-rank restrictions are separate.
This proves the stated conditional independence.
For example, before imposing injectivity of \([P,X]\), any dependence with a nonzero \(X\)-coefficient has probability at most \(2^{-(N-d)}\); union over its at most \(2^{d+h}\) coefficient choices gives a bound \(2^{2d+h-N}\).
The estimate for \(Y\) is identical.

The marginal of \(X\) alone is invariant under \(\GL(\F^N)\) and supported on injective maps; transitivity gives its uniformity, and likewise for \(Y\).
Before conditioning independent channel columns to be injective, their evaluations on a fixed independent \(q\)-tuple form a uniform \(h\)-by-\(q\) matrix.
For each nonzero coefficient vector in \(\F^q\), the chance that its evaluation is zero is \(2^{-h}\).
Union over those coefficient vectors gives at most \(2^{q-h}\), and channel injectivity has probability at least \(1/2\) for large \(N\).
This proves \Cref{eq:channel-transpose-failure}.
To compare their joint marginal with completely independent columns, condition on \(G=X\trans Y\).
The conditional marginal is uniform on the orbit of two individually injective \(h\)-frames having Gram matrix \(G\).
Under independent uniform columns that orbit has probability at least \(2^{-h^2-2}\) for large \(N\), by the argument for (d) below.
Since the raw probability of any \(G\) is at most one, the density ratio on every such orbit is at most \(2^{h^2+2}\).
Multiply over the fixed number of components to obtain the joint channel bound.
Its possible dependence on \(h\) is harmless: \(h\) is fixed before \(n\) grows.

We next prove (b), keeping explicit the dependence on the rank.
Replace each tested tuple by a basis of its nominal span; inconsistent dependent prescriptions have probability zero.
For the plus directions at the two endpoints, start with completely independent uniform plus columns and then impose individual full-frame injectivity.
The image of a rank-\(t_+\) tuple in the joint nominal direct sum is exactly uniform on \(t_+\) independent ambient vectors before this conditioning.
Indeed, at each ambient coordinate the coefficient map is a surjective linear map onto \(\F^{t_+}\), and different ambient coordinates use independent input bits.
Full injectivity has probability \(1-2^{-\Omega(N)}\), uniformly over the given construction, so the plus prescription costs at most
\((1-2^{-\Omega(N)})^{-1}2^{-t_+N}\).

Now condition on all plus frames.
At a component let \(S\) be the span of both endpoints' plus columns; then \(\dim S\le2(d+h)\).
Before minus injectivity, every minus column has an affine distribution whose translation space contains \(S^\perp\).
Split off an independent uniform \(S^\perp\)-part of each column and condition on the remaining parts.
For a rank-\(t_-\) tuple of minus coefficient combinations, the \(S^\perp\)-parts of its images are independent uniform vectors in \(S^\perp\), because the coefficient map has rank \(t_-\).
Thus prescribed minus images have probability at most
\[
 2^{-(N-2(d+h))t_-}.
\]
The additional conditioning on the two minus frames being injective changes this by a factor \(1+2^{-\Omega(N)}\), by (c).
Perform this argument componentwise.
There is a fixed total conditioning factor bounded, for example, by \(4^{|\cE|}\), and the resulting bound is
\[
 4^{|\cE|}\,2^{-t_+N-(N-2(d+h))t_-}.
\]
Choose \(M_0\) so that \(2(d+h)/N<\epsilon/2\) for large \(n\), and then increase \(n\) so that the fixed prefactor is at most \(2^{\epsilon N/2}\).
Since \(t=t_++t_-\ge1\), this proves \Cref{eq:reference-image-cap} uniformly for every allowable \(t\).
In particular, mixing the two endpoints in a direction causes no loss of an entire \(N\)-bit image constraint.

Finally, to prove (d), the plus columns are independent with probability at least \(1-2^{p-N}\): each nonzero coefficient dependence has probability \(2^{-N}\).
Given their independence, each minus column takes its required evaluations with probability \(2^{-p}\), independently, so the Gram costs \(2^{-pq}\).
Under this Gram conditioning, every nonzero combination of minus columns is uniform on an affine space of dimension at least \(N-p\).
Its probability of vanishing is at most \(2^{-(N-p)}\).
Union over fewer than \(2^q\) combinations gives failure probability at most \(2^{p+q-N}\).
This proves the lower bound.
For the upper bound with plus injectivity required, condition on the plus columns and use the exact \(2^{-pq}\) Gram probability whenever they are independent.
A partially specified consistent Gram contains some full specification; for bounded \(p,q\), the same positive lower bound for that full specification proves the last assertion.
\end{proof}

\subsection{Thinning the channel images}

Write \(e_* = |\cE|\), \(d=\dim\cB\), and
\[
 q_h=2^h-1,\qquad k_h=2e_*q_h,\qquad
 \varepsilon_{\rm ch}=10e_*2^h\rho.
\]
After the channel dimension \(h\) has been fixed, choose \(\rho>0\) so small
that \(10^4e_*2^{2h}\rho<\zeta/100\) and that this quantity is smaller than
every fixed rate slack used below.  We then choose \(M_0\) large enough
also in terms of \(1/\rho\).  The channel frame in sign \(+\) is \(X_e\),
and that in sign \(-\) is \(Y_e\).

\begin{lemma}[Channel thinning]\label{lem:channel-thinning}
Fix any bound \(q_{\rm test}\) on the lengths of the ambient tuples to be tested;
this bound is fixed before \(\rho,M_0,n\).  For all sufficiently large \(n\)
there are subsets \(A_e^\pm\subseteq V_e^\pm\), each of density
\((1+o(1))2^{-\rho N}\), with the following properties.
Keep the \(\rawlaw\)-marginal on all primal frames \((P_e,Q_e)_e\).
Given those frames, condition the original channel law on
\begin{equation}\label{eq:channel-allowed}
 X_e c\in A_e^+,\qquad Y_e c\in A_e^-
 \quad(e\in\cE, 0\ne c\in\F^h).
\end{equation}
This conditional event has positive probability for every primal frame.
The resulting probability law \(\pi\) satisfies:
\begin{enumerate}[label=\textup{(\alph*)}]
\item Its primal marginal is exactly that of \(\rawlaw\), and
\begin{equation}\label{eq:thinning-density}
 \pi^2\le 2^{\varepsilon_{\rm ch}N}\rawlaw^2.
\end{equation}
\item For every fixed independent ambient \(q\)-tuple, \(q\le q_{\rm test}\),
transpose evaluation by any specified channel frame has failure probability
at most \(2^{q-h+1}+o(1)\) under \(\pi\), uniformly in the tuple.
More generally every fixed finite union of these rank-failure tests has its
\(\rawlaw\)-probability plus \(o(1)\) as an upper bound.
\item Given any primal frames, any component and sign, and any nonzero
\(v\in V_e^\pm\),
\begin{equation}\label{eq:channel-translate}
 \pi\bigl\{\exists c\ne0:\ X_ec\in A_e^++v
              \mid(P,Q)\bigr\}\le2^{-\rho N/2},
\end{equation}
with the analogous statement for \(Y_e,A_e^-\).
\end{enumerate}
The error in \textup{(b)} is exponentially small in \(N\).
The ambient tuple there is fixed before sampling the vertex; no uniform
conditional rank bound for a tuple selected from the sampled primal frame
is asserted.
\end{lemma}
\begin{proof}
Include every ambient vector in each \(A_e^\pm\) independently with
probability \(p=2^{-\rho N}\); subsets for different components and signs
are independent.  Fix a complete primal frame \(f\), and let \(\nu_f\)
be the original conditional channel law.  Define its survival fraction
\(Z_f=\nu_f(\text{\Cref{eq:channel-allowed}})\), a random variable in the
random subsets.  Channel injectivity makes the \(q_h\) nonzero values in
each component and sign distinct.  Consequently
\[
 \mathbb E_A Z_f=p^{k_h}.
\]
For any transpose-rank test \(T\), or any union of at most \(N\)
such tests, the numerator
\(Z_{f,T}=\nu_f(T\cap\text{\Cref{eq:channel-allowed}})\) similarly has
expectation \(p^{k_h}\nu_f(T)\).

We establish simultaneous concentration of these fractions by the
bounded-differences method \citep{azuma1967,mcdiarmid1989}, with the
needed exponential-moment estimate included below.  By
\Cref{lem:frame-images}(c), a prescribed nonzero channel combination,
conditional on \(f\), has probability at most \(2^{-(N-d)+1}\) of taking
any prescribed value for all sufficiently large \(N\).  Indeed before
full-frame injectivity it is uniform on the appropriate annihilator,
and injectivity has probability \(1-2^{-\Omega(N)}\), uniformly in \(f\).
Changing one membership bit of a subset can therefore change any such
fraction by at most
\[
 c_N=C_h2^{-N+d},
\]
where \(C_h\) is fixed in \(n\).  There are \(2e_*2^N\) membership bits,
so the sum of their squared influence bounds is at most
\[
 S_N=C_{g,h}2^{-N+2d}.
\]
For clarity, expose these independent bits and use the Doob martingale
of any one fraction.  A centered increment whose conditional range has
length at most \(c\) has exponential moment at most \(\exp(t^2c^2/8)\).
To see this, the second derivative of its log moment generating function
is the variance under an exponential tilt, at most \(c^2/4\); integrate
twice using value and derivative zero at the origin.  Multiplication of
the conditional bounds and Markov's inequality give
\[
 \Pr_A(|Z-\mathbb E_AZ|>u)\le2\exp(-2u^2/S_N).
\]
Take \(u=\epsilon_N=2^{-N/10}\).  Since \(d/N\) can be made arbitrarily
small by increasing \(M_0\), these tails are at most
\(2\exp(-c_{g,h}2^{N/2})\).  There are \(2^{O(N^2)}\) primal frames and \(2^{O(N)}\) primitive
transpose tests, since their tuple lengths are bounded by the fixed
\(q_{\rm test}\).  The unions of at most \(N\) primitive tests therefore
number at most \(2^{O(N^2)}\) as well.  A union bound gives simultaneous
error at most \(\epsilon_N\) for all these denominators and test numerators.
The same argument for \(2^{-N}|A_e^\pm|\), whose bit influences are
\(2^{-N}\), gives \(|A_e^\pm|/2^N=(1+o(1))p\).

We include the translated numerators in this simultaneous estimate.
Fix \(f,e,\pm,v\ne0\), and let \(Z_{f,v}^{\pm,e}\) be the original
conditional probability of survival and of the event in
\Cref{eq:channel-translate}.  Consider, for example, sign \(+\).
If \(X_ec+v\) coincides with another already required nonzero channel
value \(X_ec'\), then \(v=X_e(c+c')\) is a nonzero channel value.
The original probability of this possibility is at most
\(C_h2^{-N+d}\), by the point bound and a union over the nonzero
coefficients.  On that exceptional event we bound the extra test by one;
the survival probability over the random subsets is still \(p^{k_h}\).
In every other case \(X_ec+v\) is a new membership bit and supplies an
additional factor \(p\).  This includes the case \(X_ec+v=0\), since
zero is not among the required nonzero channel values.  Union over \(c\)
yields
\[
 \mathbb E_AZ_{f,v}^{+,e}
 \le p^{k_h}\bigl(C_h2^{-N+d}+q_hp\bigr).
\]
A changed membership bit can affect this numerator only by being an
original channel value or a translated channel value.  Both have point
probability at most \(C_h2^{-N+d}\); hence the same influence and
concentration estimates apply.  The additional \(2^N\) choices of
translate do not affect the union bound.  We may therefore fix one
realization of all subsets for which every displayed concentration
statement holds simultaneously.

Our choice of \(\rho\) gives \(k_h\rho<1/100\).  Thus
\[
 Z_f=p^{k_h}(1+O(\epsilon_N/p^{k_h}))>0
\]
uniformly, and the relative error tends to zero exponentially.
The pointwise density of \(\pi\) with respect to \(\rawlaw\) is the
survival indicator divided by \(Z_f\).  Consequently
\[
 \frac{d\pi^2}{d\rawlaw^2}
 \le4p^{-2k_h}\le2^{\varepsilon_{\rm ch}N}
\]
for large \(N\); the primal law was left unchanged by definition.
Division of \(Z_{f,T}\) by \(Z_f\) gives
\(\pi(T\mid f)=\nu_f(T)+O(\epsilon_N/p^{k_h})\), uniformly in every
specified fixed test.  Average over the original primal marginal and
apply \Cref{eq:channel-transpose-failure} to prove \textup{(b)}.  Every
fixed finite union is among our tested unions once \(N\) is large enough;
applying the same numerator estimate directly to that union gives its
raw probability plus the stated error, without replacing the union
probability by the sum of its constituent probabilities.  Finally,
dividing the translated
numerator by \(Z_f\) bounds it by
\[
 C_h2^{-N+d}+q_h2^{-\rho N}
       +O(\epsilon_N/p^{k_h})\le2^{-\rho N/2}
\]
for all sufficiently large \(N\).  This proves \textup{(c)} and the lemma.
\end{proof}

\subsection{The two distribution tests}

For independent units \(A=(A_1,A_2)\), \(B=(B_1,B_2)\), write
\(A\bowtie B\) when all four pairs \((A_i,B_j)\) are holes.  For a point
\(z\), write \(A\bowtie z\) when both \((A_i,z)\) are holes.
These definitions also cover repeated raw types; the absence of hole
loops automatically forbids a conflict with a shared endpoint.

\begin{theorem}[Distribution tests]\label{thm:distribution-tests}
There is a constant \(\eta_{\rm col}>0\), depending only on the early
parameters and \(M\), with the following property.  Fix any \(L<\infty\)
after \(\eta_{\rm col}\) and before the selector and channel parameters;
choose all remaining parameters in the order of
\Cref{eq:parameter-order,sec:parameters}.  For every sufficiently large
\(n\), the law \(\pi\) constructed above satisfies both conclusions.
\begin{enumerate}[label=\textup{(\roman*)}]
\item Let \((A,c)\) have a probability law on units marked by a colour in
any fixed finite palette.  The colour may be randomized conditional on
\(A\).  Write \(\sigma_1,\sigma_2\) for the endpoint
marginals of its projected unit law \(\sigma\), and suppose
\begin{equation}\label{eq:coloured-law-caps}
 \sigma_1\le M\pi,\qquad \sigma_2\le M\pi,\qquad
 \sigma\le2^{3000gN}\pi^2,
 \qquad \Pr(c=j)\le\eta_{\rm col}\quad\text{for every }j.
\end{equation}
For two independent draws \((A,c),(B,c')\),
\[
 \Pr(A\bowtie B, c\ne c')\ge2^{-100gN}.
\]
\item Let \(A\) have a unit law \(\sigma\) and let \(z\) be an independent
point with law \(\theta\), where
\begin{equation}\label{eq:asymmetric-law-caps}
 \sigma_1\le L\pi,\qquad \sigma_2\le L\pi,\qquad
 \sigma\le2^{3000gN}\pi^2,
 \qquad \theta\le L\pi.
\end{equation}
Then \(\Pr(A\bowtie z)\ge2^{-100gN}\).
\end{enumerate}
The threshold for \(n\) may depend on the fixed palette and on \(L\).
All inequalities between measures are pointwise on their finite spaces.
\end{theorem}

The two conclusions are proved after the collision and overlap arguments,
in \Cref{prop:coloured-distribution,sec:asymmetric-test}, respectively.
The finite graph transfer is given in \Cref{sec:graph-transfer}.
Here we establish their leaf and image hypotheses.  In particular,
\Cref{eq:thinning-density} places every projected test law below
\(2^{(3000g+\varepsilon_{\rm ch})N}\rawlaw^2\), leaving strict slack
below the absolute budget \(D=4000g\).

\subsection{Sparse profiles arising from true intersections}

For a matrix profile, its \emph{total component rank} means \(\sum_{e\in\cE}\rank x_e\).

\begin{lemma}\label{lem:raw-intersections}
Let \(\sigma\) satisfy the joint cap in \Cref{eq:raw-law-caps}.
Outside a set of \(\sigma\)-probability \(2^{-\Omega(N)}\),
\[
 \sum_{e\in\cE}
  \dim\bigl(\im P_{1,e}\cap\im P_{2,e}\bigr)\le2D.
\]
For any unit outside that exceptional set and any
\(U_1\lambda_1=U_2\lambda_2\), both \(\lambda_i\) have total component rank at most \(2D\).
Each has a unique representation \((w_{i,l})_l\) vanishing at a strict majority of tags, and
\[
 |\{l:w_{i,l}\ne0\}|\le4D/g,\qquad
 \chi_*(w_1)=\chi_*(w_2).
\]
\end{lemma}
\begin{proof}
If the summed intersection dimension exceeds \(2D\), choose \(2D\) independent common-image vectors, componentwise.
Write each as \(P_{1,e}a=P_{2,e}b\).
Their nominal relation directions \((a,b)\) are independent in the direct sum of the two primal coefficient spaces at that component.
Indeed a dependence between the pairs would, after applying \(P_{1,e}\), be the corresponding dependence between the selected common-image vectors.
Across components their ranks add.

There are at most \(2^{O(n)}\) ways to choose this bounded number of coefficient pairs and their component indices.
For each choice, \Cref{eq:reference-image-cap} bounds the probability of their zero images by \(2^{-(1-\epsilon)2DN}\) under \(\rawlaw^2\).
After paying \(2^{DN}\) for \(\sigma\), and making the coefficient count a sufficiently small fraction of \(N\) by the choice of \(M_0\), the union bound is \(2^{-\Omega(N)}\).

Suppose \(U_1\lambda_1=U_2\lambda_2=\xi\).
The column space of \(\xi_e\) lies in both primal plus spans.
Since multiplication by the two injective frame maps preserves matrix rank,
\[
 \rank(\lambda_{i,e})=\rank(\xi_e)
 \le\dim(\im P_{1,e}\cap\im P_{2,e}).
\]
This proves the total-rank assertion.

Take any representation \((w_l)_l\) of one of the profiles.
A component \(w_l+w_{l'}\) is nonzero whenever \(w_l\ne w_{l'}\), so there are at most \(2D\) unequal unordered pairs.
If no value occurred at a strict majority of tags, write \(c_v\) for its multiplicities.
Then
\[
 \#\{\text{unequal pairs}\}
 =\frac12\left(g^2-\sum_vc_v^2\right)
 \ge\frac12\left(g^2-\frac g2\sum_vc_v\right)
 =\frac{g^2}{4}>2D,
\]
a contradiction.
Let the majority have size \(g-z\).
There are at least \(z(g-z)\ge zg/2\) unequal pairs, hence \(z\le4D/g\).
Its common value belongs to \(W_*\) by \Cref{lem:cut-kernel}, so subtract it from every coordinate.
The resulting representation vanishes at a strict majority.
It is unique: two such representations differ by a constant, and their two majority zero sets intersect.

Finally, equality of the component tensors gives equality of their ambient traces.
By cyclicity of trace and \Cref{eq:E-moments},
\[
 \Tr(P_{i,e}\lambda_{i,e}Q_{i,e}\trans)
 =\pair{P_{i,e}\trans Q_{i,e}}{\lambda_{i,e}}
 =\eta(\lambda_{i,e}).
\]
For \(e=\{l,l'\}\), it follows that
\[
 \eta(w_{1,l})+\eta(w_{2,l})
 =\eta(w_{1,l'})+\eta(w_{2,l'}).
\]
This value is independent of the tag.
The two sparse representations together have at most \(8D/g=32000<g\) nonzero coordinates, so some tag has both entries zero.
The common value is therefore zero.
Summing over tags proves equality of the two specified \(\chi_*\)-values.
\end{proof}

\subsection{Peeling into leaves with exact pins}

For a unit \(O=(1,2)\), introduce nominal coefficient spaces at each component
\[
 \cL_{O,e}
   =\bigoplus_{i\in O}(\cB_i\oplus H_i^+),
 \qquad
 \cR_{O,e}
   =\bigoplus_{i\in O}(\cB_i\oplus H_i^-),
 \qquad H_i^\pm=\F^h.
\]
Their frame-image maps add the endpoint images in the common ambient \(V_e^\pm\).
A \emph{pin space} is a specified subspace \(D_{O,e}^\pm\) of this nominal direct sum together with specified exact images of every vector in a chosen basis.
It may mix both endpoints, and it may mix primal and channel coordinates.
Once the basis images are known, every image on the pin space is fixed.

For a tested tuple of directions, its \emph{rank modulo the pins} is the sum, over components and signs, of the dimensions of the spans of its images in the quotients by \(D_{O,e}^\pm\).
The rank is nominal; it is not computed after applying the raw frames.

\begin{lemma}[First peeling and restriction]\label{lem:peeling}
Take the constants in \Cref{eq:pin-budgets,eq:law-budgets}, and choose the reference accuracy in \Cref{eq:reference-image-cap} with \(\epsilon<\zeta/4\).
For every law \(\sigma\) satisfying the absolute joint cap \Cref{eq:raw-law-caps}, the following decompositions are available.
\begin{enumerate}[label=\textup{(\alph*)}]
\item
After discarding the exception in \Cref{lem:raw-intersections} and at most \(2^{-\zeta N}\) further mass, the law is partitioned into \emph{leaves}.
Each leaf has fixed pin spaces with total dimension at most \(K_1\), and its normalized law \(\sigma_\lambda\) satisfies
\begin{equation}\label{eq:leaf-minentropy}
 \sigma_\lambda\{\text{a specified exact-image constraint of rank }t
                 \text{ modulo the pins}\}
 \le 2^{-(1-\zeta)tN}
 \qquad(t\ge1).
\end{equation}
The normalized leaf laws satisfy a common bound \(2^{O(N)}\rawlaw^2\), with the implied constant determined by the early parameters.

\item
Suppose a restriction of the mixed law has relative mass \(2^{-o(N)}\).
One may discard an exponentially small fraction of that restricted law so that every remaining normalized leaf restriction satisfies
\begin{equation}\label{eq:final-leaf-minentropy}
 \sigma_\lambda\{\text{a specified exact-image constraint of rank }t
                 \text{ modulo the pins}\}
 \le 2^{-(1-2\zeta)tN}
 \qquad(t\ge1),
\end{equation}
and still has density \(2^{O(N)}\) relative to \(\rawlaw^2\).
For several successive restrictions, the same statement applies to their combined relative cost.

\end{enumerate}
In all parts the mixed law uses each leaf's actual mass, normalized only after the indicated discards.
No equal reweighting of leaves is used.
\end{lemma}
\begin{proof}
We give the finite procedure and its quantitative bounds.
Work first with the original, unnormalized restriction of \(\sigma\) obtained by removing the intersection exception.
In particular its density is still at most \(2^{DN}\) relative to \(\rawlaw^2\).
Let \(R\) be its residual support.
As long as its residual mass exceeds \(2^{-\zeta N}\), begin with its normalized restriction and no pins.
If \Cref{eq:leaf-minentropy} fails, condition on a violating image constraint and add its nominal directions to the pins.
Given the existing pins, a consistent tuple constraint can be reduced to a basis of its span modulo the pins.
The other prescriptions are either automatic or inconsistent; an inconsistent event cannot violate a positive probability bound.
Thus every step adds positive rank.

If the cumulative added rank is \(u\), the unnormalized mass of the current part is greater than
\[
 2^{-\zeta N-(1-\zeta)uN}.
\]
The cumulative directions are independent after taking a basis, so the absolute density cap and \Cref{eq:reference-image-cap} give the upper bound
\[
 2^{DN-(1-\epsilon)uN}.
\]
Comparing the bounds yields
\begin{equation}\label{eq:peeling-rank-comparison}
 (\zeta-\epsilon)u<D+\zeta.
\end{equation}
This is strictly below the rank budget \(K_1\).
The same comparison applies to a whole violating tuple at once, so a large tuple cannot jump beyond that budget.
The rank therefore cannot increase indefinitely.
The process reaches a part with no violating constraint; remove that entire part as a leaf and restart on the residual.
The underlying raw unit space is finite, and every leaf is nonempty, so the procedure terminates with residual mass at most \(2^{-\zeta N}\).

A leaf of rank \(u\le K_1\) has mass greater than
\(2^{-\zeta N-(1-\zeta)uN}\).
Dividing the original density cap by this lower bound shows, for example,
\[
 \sigma_\lambda\le
 2^{(D+K_1+1)N}\rawlaw^2.
\]
This proves (a), including uniformity of its density constant.

For (b), write the current normalized mixture as
\(\sum_\lambda w_\lambda\sigma_\lambda\).
Let \(q_\lambda\) be the relative mass retained by the combined restriction in leaf \(\lambda\), and let
\(q=\sum_\lambda w_\lambda q_\lambda=2^{-o(N)}\).
Discard leaves for which \(q_\lambda<2^{-\zeta N}\).
Their mass in the restricted law is at most
\[
 q^{-1}\sum_{\lambda:q_\lambda<2^{-\zeta N}}
       w_\lambda q_\lambda
 \le 2^{-\zeta N+o(N)}.
\]
On every other leaf, division by \(q_\lambda\) multiplies probabilities by at most \(2^{\zeta N}\).
For \(t\ge1\),
\[
 2^{\zeta N}2^{-(1-\zeta)tN}
 \le2^{-(1-2\zeta)tN}.
\]
This proves the claimed min-entropy and density bounds.
It also explains why the small mass bound is needed for the mixed law, rather than separately asserted for every old leaf.

Throughout these operations, disintegration over the finite parts is exact.
Multiplying each normalized part by its actual mass recovers the retained law.
This proves the last assertion as well.
\end{proof}

The per-leaf density bounds in \Cref{lem:peeling} are joint bounds relative to \(\rawlaw^2\).
They do not assert that each leaf has the original one-vertex marginal cap.
If the initial mixed law has endpoint caps \(C\pi\), restricting it by
mass \(q\) gives endpoint caps \((C/q)\pi\).  Its primal marginal caps
are therefore \((C/q)\operatorname{pr}_*\rawlaw\), where
\(\operatorname{pr}\) forgets all channel columns.  For \(q=2^{-o(N)}\)
these are \(C2^{o(N)}\) times the original \emph{primal} law.
The corresponding full-frame bound is \(C2^{o(N)}\pi\), not
\(C2^{o(N)}\rawlaw\).  The absolute joint bound is
\(2^{(D+o(1))N}\rawlaw^2\).
\Cref{sec:probability-scales} records where these distinct controls are used.

\begin{corollary}[Colours and actual leaf weights]\label{cor:coloured-peeling}
For a law satisfying \Cref{eq:coloured-law-caps}, discard colours whose
mass is below \(2^{-\zeta N}\), and apply \Cref{lem:peeling}(a) to each
remaining conditional colour law.  The total discarded mass tends to
zero exponentially, the leaves have at most \(K_1\) pins and satisfy
\Cref{eq:leaf-minentropy}, and each leaf has an absolute density bound
\(2^{(D+K_1+1)N}\rawlaw^2\).  With their actual retained masses these
leaves form a mixed law whose endpoint marginals are bounded by
\((1+o(1))M\pi\).  Any subsequent combined restriction of mass
\(2^{-o(N)}\) has the first-pin freshness in
\Cref{eq:final-leaf-minentropy}, after the trimming in
\Cref{lem:peeling}(b).
\end{corollary}
\begin{proof}
The palette size is fixed, so deleting light colours costs at most a
fixed multiple of \(2^{-\zeta N}\).  Every retained conditional colour
law has absolute density at most
\(2^{(3000g+\varepsilon_{\rm ch}+\zeta)N}\rawlaw^2\), which is below
\(2^{DN}\rawlaw^2\).  The intersection exception and the residual
mass in \Cref{lem:peeling}(a) are exponentially small uniformly over
these laws.  Multiply each normalized leaf by its actual colour mass
and its mass within that colour.  Summing recovers exactly the retained
restriction of the original law, so its endpoint marginals are bounded
by the original ones divided by the total retained mass.  This proves
the asserted pooled bounds.  The final assertion is precisely
\Cref{lem:peeling}(b), applied to this mixture.  In particular no inverse
colour mass enters the pooled marginal constant.
\end{proof}

\begin{corollary}[An independent pair as an unpinned leaf]
\label{cor:independent-point-leaf}
Let \(\theta\le L\pi\), and draw two independent \(\theta\)-points.
The probability that they form an internal hole is \(2^{-\Omega(N)}\).
Their law conditioned on being a unit is a single leaf with no pins,
with endpoint marginals at most \((1+o(1))L\pi\), and it satisfies
\Cref{eq:leaf-minentropy} for every rank \(t\ge1\).  Its effective spaces
\(C_1,C_2\), defined in \Cref{sec:moments}, are consequently zero.
\end{corollary}
\begin{proof}
Under the raw primal law, each \(P_e\) is uniform injective.  Two
independent such images have nonzero intersection with probability at
most \((2^d-1)^2/(2^N-1)\), by a union bound over nonzero coefficient
pairs.  The primal marginals under \(\theta\) have density at most \(L\)
relative to those laws.  Union over components therefore bounds the
probability of any plus-span intersection by
\(O_g(L^22^{2d-N})=2^{-\Omega(N)}\).
If the plus spans are disjoint in every component, a tensor common to
the two endpoint primal tensor spaces has every component zero: each
of its column spaces lies in the corresponding plus-span intersection.
Injectivity of the primal frame maps then makes the sharing witnesses
zero, contrary to the parity sum one in
\Cref{eq:hole-equations}.  A hole is thus included in the exceptional
event just bounded.

Before deleting internal holes the pair law is at most \(L^2\pi^2\).
For a rank-\(t\) exact-image event, \Cref{eq:reference-image-cap} and
\Cref{eq:thinning-density} bound its probability by
\[
 L^2 2^{\varepsilon_{\rm ch}N-(1-\epsilon)tN}.
\]
Take \(\epsilon<\zeta/4\).  Since
\(\varepsilon_{\rm ch}<\zeta/4\) and \(t\ge1\), division by the
nonhole probability \(1-2^{-\Omega(N)}\) leaves this at most
\(2^{-(1-\zeta)tN}\) for large \(N\), uniformly for all ranks.
The same normalization gives the endpoint caps.  Empty pin spaces have
zero primal projections, which gives the assertion about \(C_i\).
\end{proof}

\subsection{Bilinear sign estimates}

\begin{lemma}[Walsh estimates]\label{lem:walsh}
Let all weights below have absolute value at most one.
\begin{enumerate}[label=\textup{(\alph*)}]
\item
If \(x,y\) are independent uniform vectors and \(B(x,y)\) is a bilinear form of rank \(d\), then
\[
 \left|\E f(x)g(y)(-1)^{B(x,y)}\right|\le2^{-d/2}.
\]
Replacing either uniform law by a law of density at most \(C_i\) multiplies this bound by at most \(C_1C_2\).

\item
If \(\alpha,\beta\) are subprobability measures on \(\F^d\), with point masses at most \(p_1,p_2\), respectively, then
\[
 \left|\sum_{x,y}\alpha(x)\beta(y)f(x)g(y)(-1)^{x\cdot y}\right|
 \le2^{d/2}(p_1p_2)^{1/2}.
\]

\item
For independent uniform global \(N\)-bit vector slots in two groups, a cross-pairing character with coefficient matrix \(C\) has bit rank \(N\rank C\).
In particular, if its cross coefficient pattern is nonzero, its correlation against two bounded weights, one on each group, is at most \(2^{-N/2}\).
Arbitrary phases depending separately on the groups may be included in these weights.
\end{enumerate}
\end{lemma}
\begin{proof}
Let \(H\) be the \(2^d\)-square matrix with entries \(H_{xy}=(-1)^{x\cdot y}\).
For \(x,x'\in\F^d\),
\[
 (HH\trans)_{xx'}=\sum_y(-1)^{(x+x')\cdot y}
 =\begin{cases}2^d,&x=x',\\0,&x\ne x'.\end{cases}
\]
The second value is zero by pairing \(y\) with \(y+e_j\) at any coordinate where \(x+x'\) is nonzero.
Thus the Euclidean operator norm of \(H\) is \(2^{d/2}\).

For (b), set \(a_x=\alpha(x)f(x)\) and \(b_y=\beta(y)g(y)\).
Then
\[
 |a\trans Hb|\le2^{d/2}\|a\|_2\|b\|_2,\qquad
 \|a\|_2^2\le p_1\sum_x\alpha(x)\le p_1,
\]
and similarly for \(b\), proving the result.

For (a), invertible changes of coordinates reduce the form to the dot product on its \(d\) active coordinates.
Average each bounded weight over its inactive coordinates.
Apply (b) to the uniform probabilities \(p_1=p_2=2^{-d}\).
A bounded-density change can be absorbed into the weights after dividing their bounds by \(C_1,C_2\).

Finally, the matrix of the full bit form in (c) is \(C\otimes I_N\).
Choose bases reducing \(C\) to \(\operatorname{diag}(I_{\rank C},0)\); the tensor matrix then has rank \(N\rank C\).
Apply (a), absorbing any within-group phases into the corresponding bounded weight.
\end{proof}
